\chapter{导论：正交宇宙}

Greg Egan 的科幻三部曲《正交》（\emph{Orthogonal}）设定在一个时空几何与我们截然不同的宇宙中：那里的度规（metric）是黎曼型（Riemannian）的 $(+,+,+,+)$，四个维度地位完全平等，没有任何一个方向被单独挑出来充当"时间"。本章阐明这一改动的物理内涵，交代三部曲的剧情背景与本书的来龙去脉，并给出后续各章的路线图。本节及后续内容所依据的物理结论，全部来自 Egan 为小说撰写的官方物理注记（physics notes）；本书即是对这批注记的编译与整理。

\orthofig[0.55\textwidth]{images/Orthogonal200.png}{\emph{Orthogonal} 宇宙的设定：时空四维地位完全平等}

\section{一个符号之差}

我们所处的宇宙是洛伦兹型（Lorentzian）的：取光速 $c=1$，平直时空中的"勾股定理"（时空线元）写作

\begin{equation}
\dd s^2 = -\dd t^2 + \dd x^2 + \dd y^2 + \dd z^2 ,
\tag{0.1}
\end{equation}

其中时间项带着一个负号。这个负号绝非无关紧要的细节，它决定了全部因果结构：正是它把四维时空劈成类时（timelike）、类光（lightlike）与类空（spacelike）三类方向，划出光锥，赋予时间以特殊的地位。

Egan 的整个设定只改动了这一处：把这个减号换成加号，得到黎曼型（Riemannian）度规

\begin{equation}
\dd s^2 = \dd t^2 + \dd x^2 + \dd y^2 + \dd z^2 .
\tag{0.2}
\end{equation}

四维以此完全对称，不存在任何被视为"时间"的特殊方向，时空的度规处处正定，恰如一个普通的四维欧几里得空间——这一设定因此得名"正交"（orthogonal）。这一符号的翻转带来若干极具反直觉、却又自洽的后果：

\begin{itemize}
  \item \textbf{时空图可以如实画在纸上。}在洛伦兹时空中，光锥把时空图割裂开来，"同时"的截面随观察者而倾斜，凭平面几何的直觉去丈量常常出错；而在黎曼时空中，四维地位平等、度规处处正定，时空图就是一张货真价实的欧几里得图形，尺规直觉不再失效。
  \item \textbf{高速远行，对旅行者意味着更长而非更短的时间。}在洛伦兹时空中，运动时钟因时间膨胀而走得慢，远行者比留在家乡的人更年轻。黎曼时空中的情形恰好相反：一艘高速飞船的船上旅程会持续许多代人的时间，而母星上只过去了短短几年——旅行者经历的是更长的岁月。
  \item \textbf{不存在普适光速，也不存在速度上限。}由于没有特殊的时间方向，也就不存在一个所有观察者都认同的不变速度 $c$。光的传播速度不再唯一，它与光的波长（乃至产生方式）有关；相应地，也就不存在"任何速度都不能超过光速"这一上限。
  \item \textbf{光的产生会释放能量。}在我们的宇宙里，原子发光要付出能量、植物靠吸收光来获得养分；在正交宇宙中，发光过程本身是释放能量的——植物向漆黑的夜空发出自己的光，借此制造养分。
\end{itemize}

\orthofig{images/02__007.png}{恒定加速度下的世界线（World Lines of Constant Acceleration）}

这些差异会一路渗透到运动学、波、电磁学、热力学、引力乃至量子力学的每一个角落，正是本书后续各章要逐一展开的内容。

\section{三部曲与本书的由来}

《正交》由三卷构成，依次为《发条火箭》（\emph{The Clockwork Rocket}）、《永恒之火》（\emph{The Eternal Flame}）与《时间之箭》（\emph{The Arrows of Time}）。故事的主人公雅尔达（Yalda）生活在这样一个宇宙里。

开篇时雅尔达不到三岁。她把孱弱的外公达里奥（Dario）带进林中静养——达里奥连日萎靡不振，不肯离开一家人安睡的花圃。村中的莉维亚医生（Doctor Livia）诊断他患了严重的"光照不足"：农场的作物几乎是单色的，人体需要更宽的光谱；太阳光过于偏蓝、太快而让身体无法获取，田间的光又全是缓慢的红光，老人需要的恰是介于两极之间的赭黄、藤黄、藏红与金黄、翠玉与铜绿。

在这个宇宙中，\textbf{光没有普适的速度，光的产生会释放能量}；而\textbf{植物通过向黑暗的夜空发出自己的光来制造食物}。

雅尔达幼年时目睹了一连串奇异流星中的一颗，它们以空前巨大的速度闯入行星系统，被称为"伤人者"（Hurtlers）。事态很快明朗：她的世界危在旦夕，而应对"伤人者"所需的知识与技术，远远超出她所在文明已有的水平。看来唯一可行的出路是：派出一艘飞船以足够高的速度远航。这趟航行对船上的人而言要持续许多代人，可对故乡而言，仅仅过去几年——船员们将有机会发现母星急需的科学，并及时把它带回来，在灾难降临之前扭转乾坤。

《正交》讲的正是雅尔达及其后代在漫长使命的险途中求存、并努力为自己开辟有意义的生活，而毁灭的阴影始终笼罩着他们远在身后的世界的故事。全书分为三卷：

\begin{itemize}
  \item \textbf{第一卷：《发条火箭》}（\emph{The Clockwork Rocket}）
  \item \textbf{第二卷：《永恒之火》}（\emph{The Eternal Flame}）
  \item \textbf{第三卷：《时间之箭》}（\emph{The Arrows of Time}）
\end{itemize}

三部曲陆续由 Night Shade Books（旧金山/纽约）与 Orion/Gollancz（伦敦）出版：第一卷初版于 2011 年，第二卷于 2012 年，第三卷于 2013 年；此后有日文（早川书房）、意大利文（Mondadori/Urania）与俄文（Explorer Books）等译本。

围绕这一设定，作者撰写了一整套系统的物理注记：主页的"主论述"（main exposition）只需高中数学即可读懂，而标注"[进阶]"（[Extra]）的页面则深入得多，面向受过本科物理与数学训练、希望更详细钻研该主题的读者。全套注记逾八万词、配有一百余幅插图。本书即是在这些官方注记的基础上编译整理而成的中文专着，供物理学专业研究生阅读。

\section{本书结构}

本书依照循序渐进的原则编排，从最基础的运动学逐步推进到引力与量子领域。各章及其中一句话即可概括的核心结果如下：

\begin{itemize}
  \item \textbf{加与减}：从时空"勾股定理"中的加减号出发，建立洛伦兹与黎曼两种几何的基本对照。
  \item \textbf{双勾股定理}（The Dual Pythagorean Theorem）：给出一个简洁的几何证明，导出该宇宙中若干关键结果的根基。
  \item \textbf{几何与运动}（Geometry and Motion）：考察点积、无类时方向的时间、能量—动量矢量与守恒律，揭示"奇怪碰撞"等有别于相对论的碰撞规律。
  \item \textbf{几何与波}（Geometry and Waves）：讨论平面波、标量波与矢量波，说明波的叠加与偏振结构。
  \item \textbf{黎曼电磁学}（Riemannian Electromagnetism）：处理三种偏振、多普勒频移、电磁场张量，以及"波纹状库仑力"与非主流电子学。
  \item \textbf{黎曼热力学}（Riemannian Thermodynamics）：给出"比无穷热还热"的奇特结论，并推导黑体谱与理想气体性质。
  \item \textbf{黎曼广义相对论}（Riemannian General Relativity）：说明为何没有黑洞、没有引力波，并（很可能）没有大爆炸。
  \item \textbf{黎曼量子力学}（Riemannian Quantum Mechanics）：讨论固体结构、非理想谐振子、固体中的量子效应、黎曼原子与高分子。
  \item \textbf{世界观与文明设定}：梳理该宇宙中物理如何塑造其文明形态的诸般设定。
  \item \textbf{术语表}（Glossary）：汇总全书术语，便于随时查阅对照。
\end{itemize}

\orthofig{images/03__008.png}{略微不同的平面波之和（Sum of slightly different plane waves）}

按照官方注记的分类，本书同样保留了进阶部分的深入内容——这包括：双勾股定理的进阶讨论与对偶矢量；运动学中的对称性、自由粒子的哈密顿量与拉格朗日量；波论中关于对称群的表示、$SO(3)$ 与 $SO(4)$ 的双覆盖、自旋表示、四元数与李代数；电磁学的黎曼 Proca 方程、静电场与静磁场的各类算例、电磁能量流动与应力—能量张量；热力学的黑体谱与单粒子配分函数；广义相对论中的黎曼史瓦西解、引力红移与时间膨胀、宇宙学模型与引力波；量子力学中的黎曼狄拉克方程、黎曼伽马矩阵、平面波解与四元数方法；以及黎曼量子场论中的正则量子化、粒子与反粒子、传播子等内容。读者可根据需要选读。

\orthofig{images/03__009.png}{$SO(3)$ 中的不可收缩回路（Non-contractible loops in SO(3)）}

\orthofig{images/04__007.png}{点电荷周围的电通量（Electric Flux Around A Point Charge）}

\orthofig{images/05__005.png}{黎曼腔体辐射谱（Riemannian Cavity Radiation Spectrum）}

\orthofig{images/06__004.png}{黎曼引力散射（Riemannian gravitational scattering）}

\orthofig{images/07__011.png}{$SO(4)$ 不变的传播子（SO(4)-invariant propagator）}

\begin{keybox}[核心设定速记]
正交宇宙与我们宇宙的物理差异，可以浓缩为以下几点：

\begin{enumerate}
  \item \textbf{度规全正}：时空度规为黎曼型 $(+,+,+,+)$，而非洛伦兹型 $(-,+,+,+)$。
  \item \textbf{四维平等}：没有任何方向被单独指定为"时间"，不存在有特殊地位的类时方向。
  \item \textbf{时空图可如实作画}：正定的欧几里得式几何使尺规直觉在时空图中同样有效。
  \item \textbf{远行耗时而非省时}：高速星际旅行对旅行者意味着更长（许多代）的岁月，对故乡却只过数年。
  \item \textbf{无普适光速}：光速与波长（及产生方式）有关，不存在所有观察者共认的不变速度。
  \item \textbf{无速度上限}：不存在任何速度不得超越的极限。
  \item \textbf{发光释放能量}：光的产生过程释放能量，植物靠向夜空发光来制造养分。
  \item \textbf{引力与量子皆异}：没有黑洞、没有引力波，很可能也没有大爆炸；原子、固体与场论均须按黎曼几何重建。
\end{enumerate}
\end{keybox}

上面这一条条差异，都源自本章开篇处那一个符号的翻转。接下来的章节将逐一展开它们背后的数学与物理。
